% ============================================================================
% CHEATSHEET TEMPLATE
% ============================================================================
% Based on the ANA-G cheatsheet style
%
% Features:
% - 4-column landscape layout on A4
% - Colored boxes for main concepts and sub-concepts
% - Compact formatting optimized for information density
% - Mathematical typesetting support
%
% Usage:
% 1. Replace metadata (title, author, date)
% 2. Adjust number of columns if needed (see \Ncolumn)
% 3. Add your content using mainbox and subbox environments
% 4. Compile with pdflatex (run twice for proper layout)
% ============================================================================

\documentclass[a4paper, fontsize=7pt]{scrartcl}

% ============================================================================
% ENCODING & LANGUAGE
% ============================================================================
\usepackage[utf8]{inputenc}
% \usepackage[german]{babel}    % Uncomment for German hyphenation
% \usepackage[english]{babel}   % Uncomment for English hyphenation

% ============================================================================
% PAGE LAYOUT
% ============================================================================
% 4-column landscape layout with minimal margins
\usepackage[landscape, left=0.75cm, top=1cm, right=0.75cm, bottom=1.5cm, footskip=5pt]{geometry}
\usepackage{flowfram}
\ffvadjustfalse
\setlength{\columnsep}{1cm}

% Number of columns (adjust as needed: 3-4 recommended for A4)
\Ncolumn{4}

% Remove paragraph indentation
\setlength{\parindent}{0cm}

% ============================================================================
% COLORS & BOXES
% ============================================================================
\usepackage[most]{tcolorbox}

% Base style for all boxes
\tcbset{
  base/.style={
    boxrule=0.3mm,
    leftrule=1mm,
    left=1.75mm,
    arc=0mm,
    fonttitle=\bfseries,
    colbacktitle=black!10!white,
    coltitle=black,
    toptitle=0.75mm,
    bottomtitle=0.25mm,
    title={#1}
  }
}

% Main concept box (green accent)
% Use for: definitions, theorems, important formulas
\definecolor{maincolor}{rgb}{0.004, 0.5, 0.35}  % Nice green
\newtcolorbox{mainbox}[1]{
  colframe=maincolor,
  base={#1}
}

% Sub concept box (gray accent)
% Use for: examples, remarks, minor concepts
\newtcolorbox{subbox}[1]{
  colframe=black!20!white,
  base={#1}
}

% Alternative colors (uncomment to use)
% \definecolor{maincolor}{rgb}{0.2, 0.4, 0.8}   % Blue
% \definecolor{maincolor}{rgb}{0.8, 0.2, 0.2}   % Red
% \definecolor{maincolor}{rgb}{0.5, 0.2, 0.6}   % Purple

% ============================================================================
% MATHEMATICS
% ============================================================================
\usepackage{amsmath, amsthm, mathtools, amssymb}
\usepackage{amsfonts, amstext}
\allowdisplaybreaks  % Allow page breaks in equations

% Common math shortcuts (customize as needed)
\def\limn{\lim_{n\to \infty}}
\def\limxo{\lim_{x\to 0}}
\def\limxi{\lim_{x\to\infty}}
\def\limxn{\lim_{x\to-\infty}}
\def\sumk{\sum_{k=1}^\infty}
\def\sumn{\sum_{n=0}^\infty}

% Number sets
\def\N{\mathbb{N}}  % Natural numbers
\def\Z{\mathbb{Z}}  % Integers
\def\Q{\mathbb{Q}}  % Rationals
\def\R{\mathbb{R}}  % Reals
\def\C{\mathbb{C}}  % Complex

% Calculus
\def\dx{\,dx}
\def\dt{\,dt}
\def\dy{\,dy}
\newcommand{\ddx}[1]{\frac{d#1}{dx}}
\newcommand{\ddn}[2]{\frac{d#1}{d#2}}
\newcommand{\pdn}[2]{\frac{\partial #1}{\partial #2}}

% Big O notation
\def\oof{\mathcal{O}}

% ============================================================================
% TABLES
% ============================================================================
\usepackage{tabularx, multirow, booktabs}
\renewcommand*{\arraystretch}{1.5}  % More spacing in tables

% ============================================================================
% LISTS
% ============================================================================
\usepackage{enumitem}
\setitemize{itemsep=0.5pt, topsep=2pt}
\setenumerate{itemsep=0.75pt, topsep=2pt}

% ============================================================================
% GRAPHICS & TIKZ
% ============================================================================
\usepackage{graphicx}
\usepackage{tikz}
\usetikzlibrary{
  positioning,
  calc,
  fit,
  matrix,
  decorations,
  angles,
  quotes,
  babel,
  3d,
  arrows.meta,
  shadings,
  shapes.geometric
}

\tikzset{
  >=stealth,
  % Add custom TikZ styles here
}

% ============================================================================
% HYPERLINKS
% ============================================================================
\usepackage{hyperref}
\hypersetup{
  colorlinks=true,
  linkcolor=maincolor,
  urlcolor=maincolor
}

% ============================================================================
% ADDITIONAL PACKAGES
% ============================================================================
\usepackage{centernot}   % For \centernot command
\usepackage{marvosym}    % Additional symbols
\usepackage{wasysym}     % Additional symbols

% ============================================================================
% DOCUMENT METADATA
% ============================================================================
\title{Cheatsheet ITEO}
\author{Thomas Schnyder}
\date{September 2026}

% Table of contents depth
\setcounter{tocdepth}{4}
\setcounter{secnumdepth}{4}

% ============================================================================
% DOCUMENT START
% ============================================================================
\begin{document}

\maketitle

% ============================================================================
% SECTION 1: DAS DATA CENTER
% ============================================================================
\section{Das Data Center}

\subsection{Building Blocks}

\begin{mainbox}{Building Blocks eines RZ}
\begin{itemize}
    \item \textbf{Gebäude} (Standort, bauliche Massnahmen)
    \item \textbf{Zutrittskontrolle}
    \item \textbf{Brandschutz / Hochwasserschutz}
    \item \textbf{Klimatisierung} (Kühlung, Feuchte, Luftreinheit)
    \item \textbf{Serverracks} mit Rechnern
    \item \textbf{Stromversorgung} (USV, Generator, Verteilung)
    \item \textbf{Datenleitungen / Netzwerk}
    \item \textbf{Administration \& Überwachung}
\end{itemize}
\end{mainbox}

\begin{mainbox}{Klimatisierung}
\textbf{Zweck:} Abfuhr von Wärme -- \emph{jeder} Betrieb elektronischer
Komponenten erzeugt Abwärme.

\textbf{Konstante Luftparameter:}
\begin{itemize}
    \item Temperatur: ca. \textbf{26\,°C} als Optimum empfohlen
    \item Rel. Luftfeuchte: \textbf{40--50\,\%}
    \begin{itemize}
        \item zu feucht $\Rightarrow$ Kondenswasser
        \item zu trocken $\Rightarrow$ ESD (elektrostat. Entladungen)
    \end{itemize}
    \item Reine Luft (Staub, Pollen, Insekten): verstopfte Lüfter/Kühlkörper
          $\Rightarrow$ reduzierte Kühlleistung $\Rightarrow$ Ausfall
\end{itemize}
\textbf{Probleme:} Redundanz einplanen, Kondenswasser am Kühlaggregat
(Abläufe, Pumpen), Filterkontrolle via \emph{Differenzdruckmessung}, hoher
Energiebedarf, Luftverteilung sicherstellen, Überwachung 7$\times$24\,h.

\textbf{Kühlleistung} oft in BTU/h:
$$\boxed{1000\ \text{BTU/h} \approx 293\ \text{W}}$$
BTU = Energie, um 1 lb Wasser um 1\,°F zu erwärmen.
\end{mainbox}

\begin{subbox}{Kühlluftverteilung}
\renewcommand*{\arraystretch}{1.2}
\begin{tabularx}{\linewidth}{@{}X|X@{}}
\textbf{Free Flow} & \textbf{Gang-Einhausung} \\
\hline
+ einfacher Aufbau & + Trennung Kalt-/Warmluft \\
+ keine zus. Hindernisse & + geringere Energiekosten \\
-- Kalt-/Warmluft mischt sich & + kleine vertikale Temp.-Diff. \\
-- ungleiche Temperaturen & -- Mehrkosten im Aufbau \\
 & -- erhöhter Abschottungsbedarf \\
\end{tabularx}
\end{subbox}

\begin{subbox}{Alternative Kühlung}
\begin{itemize}
    \item \textbf{Wasserkühlung:} nur im Geräteinnern (AIO) oder Wärmetransport
          per Wasser für ganze Racks (Wasseranschluss am Server)
    \item \textbf{Tauchkühlung} (Immersion Cooling, z.B. 3M): Hardware in
          nicht-leitender Flüssigkeit
\end{itemize}
\end{subbox}

\begin{mainbox}{USV (UPS) -- unterbrechungsfreie Stromversorgung}
Überbrückt kurze Versorgungsunterbrüche bis zum Einspringen des Generators.
\begin{itemize}
    \item \textbf{Online-USV:} Dauerbetrieb, \emph{kein} Unterbruch, aber höhere
          Verluste
    \item \textbf{Standby-USV:} sehr kurzer Unterbruch (1--20\,ms)
\end{itemize}
\begin{itemize}
    \item Deckt oft nur \emph{einen} der redundanten Strompfade ab
          $\Rightarrow$ Geräte mit 2 Netzteilen
    \item Dauer abhängig von Batteriegrösse (V, Ah)
    \item Ohne Generator: nach definierter Zeit \emph{kontrolliertes Herunterfahren}
    \item Kleine Anlagen: Batterien stark beansprucht $\Rightarrow$ Ersatz nach
          langem Ausfall
    \item Batterielebensdauer: \textbf{7--12 Jahre}
\end{itemize}
\end{mainbox}

\begin{mainbox}{Notstrom-Generator}
Für längere Überbrückung, meist \textbf{Diesel}motor-Generator.
\begin{itemize}
    \item Startzeit \textbf{5\,s bis wenige Minuten} (kalt/warm/betriebswarm)
    \item Dieseltank meist für $>$24\,h bis mehrere Tage
    \item Braucht Luftzufuhr, Abgasanlage; Lärm; Platz \& Raumkühlung
    \item Aufwändige Integration mit USV \& Netz, Serviceintervalle
\end{itemize}
\textbf{Kette:} Netz $\to$ USV (überbrückt Sekunden) $\to$ Generator (Stunden/Tage)
\end{mainbox}

\begin{mainbox}{EDV-Einbau: Serverracks}
\begin{itemize}
    \item Standard \textbf{19''} (48.26\,cm) Breite der Einbauten
    \item $\boxed{1\ \text{HE} = 1\ \text{U} = 1.75'' = 44.45\ \text{mm}}$
    \item Rack: Höhe 21--49\,U, Tiefe 0.8--1.2\,m, Breite 0.6--1\,m
    \item Integrierte Kühlung möglich; Zuleitungen oben/unten/seitlich
\end{itemize}
Dazu: Netzwerk (Kupfer, Glas), Stromverteilung (Hohlboden oder Deckentrasse),
Klimageräte, USV, Batterieschränke, Kühlleitungen, Überwachungsgeräte.
\end{mainbox}

\subsection{Kritische Punkte}

\begin{mainbox}{Gefährdungen und Abhilfe}
\renewcommand*{\arraystretch}{1.2}
\begin{tabularx}{\linewidth}{@{}>{\raggedright\arraybackslash}p{0.36\linewidth}|>{\raggedright\arraybackslash}X@{}}
\textbf{Gefährdung} & \textbf{Abhilfe} \\
\hline
Einbruch, Diebstahl, Vandalismus, Sturm, Trümmer & bauliche Massnahmen \\
Fremdzugriff & Zutrittskontrolle, Abhörsicherheit, Firewall \\
Feuer / Rauch & Brandfrüherkennung, Löschung, Abschottung \\
Netzausfälle/-störungen & Netzfilter, USV, Generatoren \\
EM-Störfelder & Abschirmung \\
Staub, Schmutz, Wasser & Filter, Schleusen, Standortwahl, Abschottung, Pumpen \\
\end{tabularx}
\textbf{Achtung:} Klimaanlage verteilt Brandgase im ganzen RZ!
\end{mainbox}

\begin{subbox}{Überwachungsgebiete}
\begin{itemize}
    \item \textbf{Gebäude:} Türen offen/zu/verschlossen, Kameras,
          Bewegungsmelder, Zutritte
    \item \textbf{Räume:} Temperatur, Feuchte, Bewegung, Rauch, Brand, Leckage
    \item \textbf{Energie:} Netzausfall, Strom, Spannung, Leistung (kW),
          Leistungsfaktor ($\cos\varphi$)
    \item \textbf{Einzelgeräte:} Schalterstellungen, Verbrauch pro Bereich,
          Sicherungsausfall, Kurzschluss, Überlast
    \item \textbf{Generator:} Kraftstoff, Bereitschaft, Temperatur, Überlast
    \item \textbf{Klimageräte:} Temperatur, Feuchte, Filterwiderstand, Störungen
    \item \textbf{USV:} Normal-/Batterie-/Bypass-Betrieb, Ladezustand,
          Batterietemperatur
    \item \textbf{Brandmelde-/Löschanlage:} ausgelöst, Übertragung aus,
          Störung, Service
\end{itemize}
\end{subbox}

\begin{mainbox}{PUE -- Power Usage Effectiveness}
Massstab für die Energieeffizienz eines RZ (``Green IT''):
$$\boxed{\text{PUE} = \frac{\text{Gesamtenergie RZ}}{\text{Energie IT-Geräte}}}$$
\begin{itemize}
    \item \textbf{1.0:} optimal -- keine Energie für Kühlung etc.
          (theoretisches Minimum)
    \item \textbf{1.2:} guter Wert, normale moderne RZ
    \item \textbf{$>$1.4:} Optimierungsbedarf
\end{itemize}
Annähernd 1 z.B. durch freie Kühlung in kaltem Klima oder Unterwasser-RZ
(MS Project Natick).
\end{mainbox}

\begin{subbox}{Beispiel PUE}
IT-Last 500\,kW, Gesamtverbrauch 650\,kW
$\Rightarrow$ PUE $= 650/500 = 1.3$
$\Rightarrow$ 150\,kW (30\,\% der IT-Last) für Kühlung, USV-Verluste, Licht.
\end{subbox}

\subsection{Verfügbarkeitsverbesserungen}

\begin{subbox}{Kosten für Downtime (pro Stunde, Richtwerte)}
\renewcommand*{\arraystretch}{1.1}
\begin{tabularx}{\linewidth}{@{}X r@{}}
Fertigung & \$28'000 \\
Logistik, Einzelhandel & \$90'000 \\
Home-Shopping & \$113'000 \\
Medien (pay per view) & \$1.1\,Mio \\
Bank (RZ), Kreditkarten & \$2.5--2.6\,Mio \\
Broker & \$6.5\,Mio \\
\end{tabularx}
\end{subbox}

\begin{mainbox}{Das Rätsel der Neunen}
$$\boxed{\text{Downtime} = (1 - A) \cdot \text{Betriebszeit}}$$
Betriebszeit/Jahr: \textbf{7$\times$24:} $365 \cdot 24 = 8760$\,h \quad
\textbf{5$\times$9:} $261 \cdot 9 = 2349$\,h
\renewcommand*{\arraystretch}{1.15}
\begin{center}
\begin{tabular}{r|r|r}
\textbf{Uptime} & \textbf{7$\times$24} & \textbf{5$\times$9} \\
\hline
90\,\% & 876\,h (36.5\,d) & 234.9\,h \\
95\,\% & 438\,h (18.25\,d) & 117.5\,h \\
99\,\% & 87.6\,h (3.65\,d) & 23.5\,h \\
99.9\,\% & 8.76\,h & 2.35\,h \\
99.99\,\% & 52.56\,min & 14.1\,min \\
99.999\,\% & 5.26\,min & 1.41\,min \\
99.9999\,\% & 31.5\,s & 8.5\,s \\
\end{tabular}
\end{center}
\textbf{Merkregel:} jede zusätzliche 9 $\Rightarrow$ Downtime $/10$
(und Kosten steigen stark).
\end{mainbox}

\begin{subbox}{Beispiel: max. 1\,h Ausfall/Jahr in Bürozeit}
Bürozeit 5$\times$9 = 2349\,h/Jahr:
$$A \geq 1 - \frac{1}{2349} = 99.957\,\% \approx \textbf{99.96\,\%}$$
Bei 7$\times$24: $1 - 1/8760 = 99.989\,\%$.
\end{subbox}

\begin{mainbox}{Massnahmen für höhere Verfügbarkeit}
\textbf{SAN spiegeln:} kostet \textbf{mehr als 100\,\%} zusätzlich -- neben dem
doppelten Speicher braucht es Replikationsleitungen, zusätzliche Ports,
Failover-Switches, Lizenzen, Platz, Strom, Betrieb.

\textbf{Failover Cluster:} Zugriff über eine \textbf{virtuelle IP}, die zusammen
mit dem Service (DB, Mail, Fileshare, Print, \dots) bei Ausfall eines Hosts
\emph{automatisch} auf den zweiten Host verschoben wird. Auch $>$2 Server mit
automatischem Lastausgleich möglich.

\textbf{Failover Datacenter:} zweites, räumlich getrenntes RZ.
\begin{itemize}
    \item \textbf{synchron:} Fibre Channel, 10G Ethernet (kurze Distanz)
    \item \textbf{asynchron:} über Internet (TCP/IP), weite Distanz
\end{itemize}
\end{mainbox}

\begin{mainbox}{Synchrone vs. asynchrone Replikation}
\begin{itemize}
    \item \textbf{Synchron:} Schreibvorgang erst bestätigt, wenn auch Kopie
          geschrieben $\Rightarrow$ kein Datenverlust, aber Latenz bremst
          primäre DB
    \item \textbf{Asynchron:} Kopie wird nachgeliefert $\Rightarrow$ schnell,
          weite Distanz möglich, aber Datenverlust beim Ausfall möglich
\end{itemize}
\textbf{Praxis -- Kombination:} schnell \emph{synchron} auf lokale Instanz
(physisch getrennt), danach \emph{asynchron} ins entfernte RZ (Standby).
Failover erst möglich, wenn Sync zu 100\,\% abgeschlossen.
\end{mainbox}

\begin{mainbox}{Availability Environment Classification (HRG)}
\renewcommand*{\arraystretch}{1.15}
\begin{tabularx}{\linewidth}{@{}l>{\raggedright\arraybackslash}X@{}}
\textbf{AEC-0} Conventional & Unterbruch ok, Datenintegrität nicht essentiell \\
\textbf{AEC-1} Highly Reliable & Unterbruch ok, Datenintegrität muss gewährleistet sein \\
\textbf{AEC-2} High Availab. & in festgelegten Zeiten nur minimal unterbrochen \\
\textbf{AEC-3} Fault Resilient & in festgelegten Zeiten \emph{ununterbrochen} \\
\textbf{AEC-4} Fault Tolerant & ununterbrochen, 24/7 \\
\textbf{AEC-5} Disaster Tol. & unter \emph{allen} Umständen verfügbar \\
\end{tabularx}
\end{mainbox}

\subsection{Tier Levels im Datacenter}

\begin{mainbox}{Data Center Tiers (Uptime Institute)}
\renewcommand*{\arraystretch}{1.15}
\begin{center}
\setlength{\tabcolsep}{4pt}
\begin{tabular}{c|l|r|r}
\textbf{Tier} & \textbf{Redund.} & \textbf{h/Jahr} & \textbf{Verf.} \\
\hline
I & N & 28.8\,h & 99.67\,\% \\
II & N+1 & 22\,h & 99.75\,\% \\
III & N+1 & 1.6\,h & 99.98\,\% \\
IV & 2$\times$(N+1) & 0.8\,h & 99.99\,\% \\
\end{tabular}
\end{center}
\begin{itemize}
    \item \textbf{N:} für Bedarf N stehen genau N Ressourcen bereit
    \item \textbf{N+1:} eine Reserve -- Ausfall einer Einheit ohne
          Leistungseinbusse
    \item \textbf{2N:} alles doppelt (vollständig redundant)
\end{itemize}
Betrifft nur die RZ-\emph{Infrastruktur} (ohne Server, Storage, Netzwerk).
\end{mainbox}

\begin{subbox}{Tiers im Detail (Auszug)}
\begin{itemize}
    \item \textbf{ab II:} Dual Utility Feed, UPS N, Chiller N+1,
          Automation \& Monitoring
    \item \textbf{ab III:} UPS N+1, Generatoren, Dual Feed CAC,
          Dual Power Path für Dual-Cord-Lasten, 7$\times$24 Operators
    \item \textbf{ab 3+:} Dual Power Path zur PDU, Generatoren N+1
    \item \textbf{IV:} UPS \textbf{2N}, Chiller N+2,
          \textbf{keine Single Points of Failure}, Concurrent Maintenance
          (Wartung im laufenden Betrieb) mech. \& elektr.
\end{itemize}
PDU = Power Distribution Unit, CAC = Central Air Conditioner.
``Tier 5'' ist kein Uptime-Standard (Marketing-Begriff einzelner Anbieter).
\end{subbox}

\subsection{Information Lifecycle Management}

\begin{mainbox}{ILM -- Information Lifecycle Management}
\textbf{Speicherstrategie:} Informationen entsprechend ihrem \textbf{Wert} auf
dem jeweils \textbf{günstigsten} geeigneten Speichermedium ablegen -- inkl.
Regeln und Prozessen, wie sie zwischen Medien verschoben werden.

Steuerung nach Wichtigkeit, Wertigkeit, Kosten. Basis: Klassifizierung von
Daten, Quellen und Speichersystemen $\Rightarrow$ automatisierte Bereitstellung
in einer Speicherhierarchie.

\textbf{Getrieben durch:} Storage Mgmt, Document Lifecycle Mgmt, Content
Lifecycle Mgmt, \textbf{Records Management}.
\end{mainbox}

\begin{mainbox}{ILM-Regeln: wo/wie wird gespeichert?}
Änderungshäufigkeit, Zugriffsgeschwindigkeit, Zugriffshäufigkeit, Kosten,
ökonomischer Wert, gesetzliche Bestimmungen.

$\Rightarrow$ Daten auf passende Medien verschieben und am Ende des Zyklus
\textbf{löschen}.
\end{mainbox}

\begin{subbox}{Aktive vs. inaktive Daten}
Datenmenge wächst exponentiell -- der Grossteil ist \textbf{inaktiv}.
Inaktive Daten trotzdem:
\begin{itemize}
    \item konsumieren teuren Speicherplatz
    \item werden gepflegt, gesichert, repliziert
    \item unterliegen Rechts- \& Aufbewahrungspflichten
    \item müssen im Katastrophenfall wiederhergestellt werden
\end{itemize}
$\Rightarrow$ auf günstigere Tiers verschieben / archivieren / löschen.
\end{subbox}

\begin{mainbox}{Records Management (ISO 15489)}
Steuert das Datenvolumen:
\begin{itemize}
    \item \textbf{Aufbewahrungszeiten} (gesetzlich / Branche / intern)
    \item Ablageordnung, Namensgebung, Versionen
    \item automatische Löschroutinen
\end{itemize}
Anordnen \& durchsetzen: \textbf{Geschäftsleitung} (mit Rechtsabteilung),
nicht Admins.
\end{mainbox}

\subsection{Tiered Storage}

\begin{mainbox}{Tiered Storage}
``Verschiedene Datenklassen auf verschiedenen Speicherklassen.''
Klassifizierung + passendes Medium \textbf{senkt die Kosten}.

\textbf{Datencharakteristik:}
\begin{enumerate}
    \item Mission Critical Data
    \item Business Critical Data
    \item Nearline / Historical Data
    \item Offline Data
\end{enumerate}
\textbf{Speicherhierarchie:}
\begin{itemize}
    \item \textbf{Tier 1:} höchster Speed, sehr zuverlässig, skalierbar,
          sehr teuer
    \item \textbf{Tier 2:} mittlerer Speed, Zuverlässigkeit ok, Skalierung
          limitiert, weniger teuer
    \item \textbf{Tier n:} hohe Kapazität, tiefer Speed, sehr geringe
          Kosten/TB
    \item \textbf{Offsite Tape:} Archivierung, WORM (einmal beschreibbar)
\end{itemize}
\end{mainbox}

\begin{mainbox}{Klassifizierung der Speicher-Tiers}
\renewcommand*{\arraystretch}{1.15}
\begin{tabularx}{\linewidth}{@{}l>{\raggedright\arraybackslash}X r@{}}
\textbf{Tier} & \textbf{Charakteristik / Gebrauch} & \textbf{Anteil} \\
\hline
Tier 1 & schnell, hochverfügbar; Produktionsdaten & 20\,\% \\
Tier 2 & mittel, hohe Dichte; Referenzdaten & 35\,\% \\
Compliance & WORM, unveränderlich; gesetzl. Daten & 30\,\% \\
Tape Vault & langsam, entfernbar; Offsite DR & 15\,\%$^*$ \\
\end{tabularx}
$^*$ plus Backups
\end{mainbox}

\begin{subbox}{Vergleich der Speichertypen}
\renewcommand*{\arraystretch}{1.15}
\setlength{\tabcolsep}{3pt}
\begin{tabularx}{\linewidth}{@{}l|l|X|l@{}}
\textbf{Typ} & \textbf{Zugriff} & \textbf{Rate} & \textbf{Kosten} \\
\hline
SSD & 0.2/0.4\,ms & 300--500\,MB/s, 20--100k IOPS & s.\,hoch \\
FC/SAS & $\sim$4\,ms & 160\,MB/s, 140 IOPS & hoch \\
SATA & $\sim$12\,ms & 100\,MB/s, 90 IOPS & moderat \\
MAID & ms--s & -- & niedrig \\
Tape & s--h & seq. steigend, rand. gering & s.\,niedrig \\
\end{tabularx}
\textbf{Dichte:} SSD/FC/SATA moderat(--hoch), MAID/Tape hoch.
\textbf{Zuverlässigkeit:} SSD (sehr) hoch, FC hoch, SATA moderat, Tape
niedrig--moderat. \textbf{Energie:} SSD/FC hoch $\to$ Tape sehr niedrig.

\textbf{MAID} = Massive Array of Idle Disks: viele Disks, stromsparend/offline,
nur bei Bedarf aktiv.
\end{subbox}

\begin{subbox}{Storage-SLA Beispiel (CS intern)}
\renewcommand*{\arraystretch}{1.15}
\setlength{\tabcolsep}{3pt}
\begin{center}
\begin{tabular}{c|r|r|r|l}
\textbf{SLA} & \textbf{Verf.} & \textbf{RTO} & \textbf{RPO} & \textbf{Kosten} \\
\hline
1 & 99.99\,\% & 15\,min & 1\,h & hoch \\
2 & 99.9\,\% & 1\,h & 12\,h & mittel \\
3 & 99\,\% & 8\,h & 48\,h & niedrig \\
4 & 97\,\% & 24\,h & 96\,h & gering \\
\end{tabular}
\end{center}
1 = Mission Critical, 2 = Business Critical, 3 = Business Important,
4 = Development.
Weitere SLA-Punkte (SLA 1$\to$4): Backup-Erfolgsrate 97/95/90/90\,\%;
Archivierung nach 90/30/90/180 Tagen ohne Zugriff; Archivzugriff Sekunden
bis 24--48\,h; Kapazitätsprognose monatlich bis jährlich; Threshold-basiertes
automatisches Nachladen von Speicher.

\textbf{RTO} (Recovery Time Objective): max. Dauer bis Wiederherstellung.
\textbf{RPO} (Recovery Point Objective): max. tolerierter Datenverlust (Zeit
seit letzter Sicherung).
\end{subbox}

\subsection{Allocation Efficiency}

\begin{mainbox}{Utilization \& Yield -- COPQ}
\textbf{COPQ} = Cost of Poor Quality: Kosten für Speicher, der bezahlt, aber
nicht genutzt wird (Reserven, RAID, Spiegelung).
\begin{align*}
\text{Yield} &= \prod \text{Effizienzfaktoren} \\
\text{COPQ} &= \text{Kosten} \cdot (1 - \text{Yield}) \\
\text{CHF/TB}_{\text{genutzt}} &= \frac{\text{Kosten}}{\text{Kap.} \cdot \text{Yield}}
\end{align*}
\textbf{Beispiel Tier 1:} 1000\,GB à \$0.05/MB $= \$51'200$;
alloziert 80\,\%, genutzt 75\,\% $\Rightarrow$ Yield $0.8 \cdot 0.75 = 60\,\%$
$\Rightarrow$ COPQ $= \$20'480$ (2.05\,\% von \$1\,Mio Einnahmen).
\end{mainbox}

\begin{subbox}{Übung: 96 Platten à 8\,TB, CHF 700.--}
\textbf{Kosten:} $96 \cdot 700 = \text{CHF } 67'200$, roh $768$\,TB.
\renewcommand*{\arraystretch}{1.15}
\begin{center}
\begin{tabular}{l|r|r}
\textbf{Schritt} & \textbf{Faktor} & \textbf{TB} \\
\hline
Roh & & 768.0 \\
Spiegelung & 0.50 & 384.0 \\
RAID-5 20\,\% & 0.80 & 307.2 \\
Reserve 20\,\% & 0.80 & 245.8 \\
Nutzung 85\,\% & 0.85 & 208.9 \\
\end{tabular}
\end{center}
\begin{align*}
\text{Yield} &= 0.5 \cdot 0.8 \cdot 0.8 \cdot 0.85 = 27.2\,\% \\
\text{COPQ} &= 67'200 \cdot 0.728 \approx \textbf{CHF 48'922} \\
\text{CHF/TB} &: 87.5 \text{ (roh)} \to 321.7
\end{align*}
\end{subbox}

\begin{subbox}{Häufige Fehler \& Tipps}
\begin{itemize}
    \item Downtime immer auf die \textbf{vereinbarte Servicezeit} (7$\times$24
          vs. 5$\times$9) beziehen, nicht blind auf 8760\,h.
    \item Spiegelung $\neq$ +100\,\% Kosten -- Infrastruktur kommt dazu.
    \item Tier I--IV (RZ-Infrastruktur) $\neq$ Storage-Tiers (Speicherklassen)
          $\neq$ AEC-Klassen (Verfügbarkeit).
    \item Effizienzfaktoren \textbf{multiplizieren}, nicht addieren.
\end{itemize}
\end{subbox}

% ============================================================================
% TIPS FOR AVOIDING OVERFULL BOXES
% ============================================================================
%
% 1. Keep formulas concise - split long equations over multiple lines
% 2. Use shorter variable names in inline math
% 3. Abbreviate table headers
% 4. Use align* environment for multi-line equations
% 5. Test compile frequently to catch overflow early
%
% Example of splitting a long formula:
% Instead of: $$a + b + c + d + e = f + g + h + i + j$$
% Use:
% \begin{align*}
% a + b + c + d + e &= f + g + h \\
%                   &\quad + i + j
% \end{align*}
%
% ============================================================================

\end{document}